%【紙面設定】
%b4,横書き,2段
\documentclass[paper=b4j,landscape,twocolumn,fleqn]{jlreq}
\special{papersize=\the\paperwidth,\the\paperheight}
\setlength{\columnseprule}{0.4pt}
\pagestyle{empty}
\usepackage[margin=8truemm]{geometry}
\usepackage[dvipdfmx]{graphicx}
\usepackage{amsmath}
\ModifyHeading{section}{lines=1,before_space=6pt,after_space=0pt}
\setlength{\parskip}{0pt}
\setlength{\jot}{3pt}

%【本文】
\begin{document}
\normalsize
\section{覚えるべき三角形の図示 解答}
\begin{align*}
&1.\ 30^\circ\text{-}60^\circ\text{-}90^\circ\text{ の辺比 }1:\sqrt{3}:2\\[0.45em]
&2.\ 45^\circ\text{-}45^\circ\text{-}90^\circ\text{ の辺比 }1:1:\sqrt{2}\\[0.45em]
&3.\ \text{最短辺}=1\Rightarrow (1,\sqrt{3},2)\\[0.45em]
&4.\ \text{斜辺}=1\Rightarrow \left(\frac{\sqrt{2}}{2},\frac{\sqrt{2}}{2},1\right)\\[0.45em]
&5.\ \sin30^\circ=\frac12,\ \cos30^\circ=\frac{\sqrt3}{2},\ \sin45^\circ=\cos45^\circ=\frac{\sqrt2}{2},\ \sin60^\circ=\frac{\sqrt3}{2},\ \cos60^\circ=\frac12\\[0.45em]
&6.\ \tan30^\circ=\frac{1}{\sqrt3},\ \tan45^\circ=1,\ \tan60^\circ=\sqrt3
\end{align*}

\section{直角三角形の比 解答}
\begin{align*}
&1.\ 2,\ 2\sqrt3,\ 4\\[0.45em]
&2.\ 4,\ 4\sqrt3,\ 8\\[0.45em]
&3.\ 6,\ 6,\ 6\sqrt2\\[0.45em]
&4.\ 30^\circ\text{側}=3,\ \text{斜辺}=6\\[0.45em]
&5.\ 60^\circ\text{側}=5\sqrt3,\ \text{斜辺}=10\\[0.45em]
&6.\ \text{もう1辺}=4,\ \text{斜辺}=4\sqrt2
\end{align*}

\section{三角比 解答}
\begin{align*}
&1.\ \cos30^\circ=\frac{\sqrt3}{2},\ \cos45^\circ=\frac{\sqrt2}{2},\ \cos60^\circ=\frac12\\[0.45em]
&2.\ \sin30^\circ=\frac12,\ \sin45^\circ=\frac{\sqrt2}{2},\ \sin60^\circ=\frac{\sqrt3}{2}\\[0.45em]
&3.\ \tan30^\circ=\frac{1}{\sqrt3},\ \tan45^\circ=1,\ \tan60^\circ=\sqrt3\\[0.45em]
&4.\ \text{底辺}=4\sqrt3,\ \text{高さ}=4\\[0.45em]
&5.\ \text{底辺}=5,\ \text{高さ}=5\sqrt3\\[0.45em]
&6.\ \text{斜辺}=6\sqrt2,\ \text{高さ}=6
\end{align*}

\newpage
\section{ラジアン 解答}
\begin{align*}
&1.\ 30^\circ=\frac\pi6,\ 45^\circ=\frac\pi4,\ 60^\circ=\frac\pi3,\ 90^\circ=\frac\pi2\\[0.45em]
&2.\ \frac\pi6=30^\circ,\ \frac\pi4=45^\circ,\ \frac\pi3=60^\circ,\ \frac\pi2=90^\circ\\[0.45em]
&3.\ \cos\frac\pi3=\frac12,\ \sin\frac\pi6=\frac12,\ \tan\frac\pi4=1\\[0.45em]
&4.\ \text{底辺}=4\sqrt3,\ \text{高さ}=4\\[0.45em]
&5.\ \text{底辺}=2\sqrt2,\ \text{高さ}=2\sqrt2\\[0.45em]
&6.\ \text{斜辺}=6,\ \text{高さ}=3\sqrt3
\end{align*}

\section{単位円の三角比 解答}
\begin{align*}
&1.\ \sin\theta=\frac{\sqrt3}{2},\ \cos\theta=\frac12,\ \tan\theta=\sqrt3\\[0.45em]
&2.\ \sin\theta=\frac{\sqrt2}{2},\ \cos\theta=\frac{\sqrt2}{2},\ \tan\theta=1\\[0.45em]
&3.\ \sin\theta=\frac{\sqrt3}{2},\ \cos\theta=-\frac12,\ \tan\theta=-\sqrt3\\[0.45em]
&4.\ \cos\theta=\pm\frac{2\sqrt2}{3},\ \tan\theta=\pm\frac{\sqrt2}{4}\\[0.45em]
&5.\ \sin\theta=\frac{\sqrt{15}}{4},\ \tan\theta=\sqrt{15}\\[0.45em]
&6.\ \sin\theta=\frac{2\sqrt5}{5},\ \cos\theta=\frac{\sqrt5}{5}
\end{align*}

\section{三角比の総合 解答}
\begin{align*}
&1.\ \frac{1+\sqrt3}{2}\\[0.45em]
&2.\ \frac12\\[0.45em]
&3.\ \frac{\sqrt3}{2}\\[0.45em]
&4.\ \sqrt3+2\\[0.45em]
&5.\ 1\\[0.45em]
&6.\ \cos\theta=\frac45,\ \tan\theta=\frac34
\end{align*}
\end{document}
